\chapter{Spectral Semantics and Curvature Measures on Polyxan Hypergraphs}
\label{ch:spectral-semantics}

Chapters~\ref{ch:polyxan-basic}--\ref{ch:polyxan-rewriting} equipped Polyxan
hypergraphs with typed structure, sheaf-valued annotations, curvature
assignments, and rewriting dynamics governed by curvature-decreasing EHR
flows.  
In this chapter we introduce the spectral geometry of Polyxan structures.
The spectral data of the normalized Laplacian, together with the sheaf
cohomology of the incidence complex, constitute a complete system of semantic
invariants.  
They are the discrete counterpart of the RSVP curvature--entropy spectrum
developed in Part~I, and they serve as the algebraic–combinatorial anchor for
all coupled field–graph dynamics in Part~III.

% ==============================================================
\section{Incidence Complex and Normalized Laplacian}
% ==============================================================

Let $G$ be a finite Polyxan hypergraph with vertex set $V$, hyperedge set $E$,
typed structure $t$, sheaf assignment $\mathscr{F}$, and curvature assignment
$\kappa_G$.

\subsection{Incidence operators}

Define the signed incidence operator
\[
\partial : C^1(G) \to C^0(G),
\qquad
(\partial f)(v) := \sum_{e \ni v} \sigma(v,e) f(e),
\]
where $\sigma(v,e) \in \{-1,0,1\}$ encodes orientation induced by
hyperedge type signatures.  
The adjoint
\[
\partial^\ast : C^0(G) \to C^1(G)
\]
is taken with respect to the curvature-weighted inner product
\[
\langle f,g\rangle_0 = \sum_{v\in V} \omega(v) f(v) g(v),
\qquad
\langle f,g\rangle_1 = \sum_{e\in E} \omega(e) f(e) g(e),
\]
where $\omega(\cdot)$ is the curvature weight induced by $\kappa_G$.

\subsection{Normalized Laplacian}

\begin{definition}[Normalized Polyxan Laplacian]
The normalized Laplacian is
\[
\Delta_G := \partial^\ast \partial,
\]
a self-adjoint, positive semidefinite operator on $C^0(G)$.
\end{definition}

Eigenfunctions of $\Delta_G$ correspond to harmonic semantic modes of the
hypergraph: they minimize curvature variation weighted by sheaf consistency.

\begin{remark}
$\Delta_G$ reduces to the standard normalized Laplacian for graphs when all
hyperedges are binary, curvature is uniform, and sheaf data is trivial.
\end{remark}

% ==============================================================
\section{Polyxan Heat Kernel and Trace Invariants}
% ==============================================================

\subsection{Heat kernel}

\begin{definition}[Heat kernel]
The heat kernel of $G$ is
\[
K_G(t) := e^{-t \Delta_G}.
\]
\end{definition}

$K_G(t)$ encodes the flow of semantic energy across the hypergraph through
diffusion shaped by curvature and sheaf constraints.

\subsection{Heat trace}

\begin{definition}[Heat trace]
The heat trace is
\[
\Theta_G(t) := \mathrm{Tr}(K_G(t))
= \sum_{\lambda \in \mathrm{Spec}(\Delta_G)} e^{-t\lambda},
\]
the generating function for the spectral density of $G$.
\end{definition}

\begin{theorem}[Spectral invariance]
$\Theta_G(t)$ is invariant under all type-preserving isomorphisms of $G$.
\end{theorem}

Thus $\Theta_G(t)$ serves as a canonical semantic invariant.

% ==============================================================
\section{Zeta Function and Analytic Torsion Analogue}
% ==============================================================

\subsection{Spectral zeta function}

\begin{definition}[Polyxan spectral zeta function]
Let $\lambda_1,\dots,\lambda_n$ be the positive eigenvalues of $\Delta_G$.
Define
\[
\zeta_G(s)
=
\sum_{i=1}^n \lambda_i^{-s},
\qquad \Re(s) > 1.
\]
\end{definition}

Analytic continuation extends $\zeta_G(s)$ to a meromorphic function on
$\mathbb{C}$.

\subsection{Analytic torsion}

\begin{definition}[Polyxan analytic torsion]
The analytic torsion analogue is
\[
T[G] := \exp\left( -\zeta'_G(0) \right).
\]
\end{definition}

$T[G]$ measures global semantic complexity:  
low torsion corresponds to high symmetry and coherent curvature flow.

% ==============================================================
\section{Spectral Action and Media-Quine Critical Points}
% ==============================================================

Define the spectral action functional
\[
S[G] := \mathrm{Tr}(\log \Delta_G)
= \sum_{\lambda>0} \log(\lambda).
\]

\begin{theorem}[Spectral critical points]
Media quines are precisely the critical points of $S[G]$ under all
curvature-preserving, type-preserving variations.
\end{theorem}

\begin{proof}[Sketch]
The EHR cycle decreases $\mathcal{E}[G]$, which corresponds to a monotone
decrease in low-frequency eigenvalues of $\Delta_G$.  
At fixed points of EHR, all curvature-lowering rewrites vanish, implying
stationarity of spectral eigenvalues.  
The differential of $S[G]$ with respect to eigenvalues is
$\delta S = \sum \delta\lambda/\lambda$, which vanishes exactly when all
$\delta\lambda = 0$.  
Thus quines ↔ spectral critical points.
\end{proof}

This parallels the role of the RSVP fixed-point Lagrangian $\mathcal{L}_\ast$ in
Chapter~\ref{ch:jet-rg}.

% ==============================================================
\section{A Polyxan c-Theorem}
% ==============================================================

Define the Polyxan central charge
\[
c_{\Polyx}[G] := \log \det \Delta_G
= \sum_{\lambda>0} \log \lambda.
\]

\begin{theorem}[Polyxan c-theorem]
Under the curvature-decreasing EHR flow,
\[
\frac{d}{dt}\, c_{\Polyx}[G_t] < 0,
\]
with equality if and only if $G_t$ is a media quine.
\end{theorem}

\begin{proof}[Sketch]
The EHR flow decreases low-frequency eigenvalues of $\Delta_G$ while leaving
null eigenspaces determined by sheaf cohomology invariant.  
Since $\log$ is strictly increasing,
\[
\frac{d}{dt} c_{\Polyx}[G_t]
= \sum_{\lambda>0} \frac{\dot{\lambda}}{\lambda}
< 0,
\]
unless all $\dot{\lambda}=0$, characterizing media quines.
\end{proof}

This is the discrete analogue of the RSVP c-theorem of
Chapter~\ref{ch:jet-rg}; it establishes the irreversibility of semantic
coarse-graining in Polyx.

\subsection{Coupled c-theorem}

\begin{theorem}[Coupled central charge unification]
At the joint fixed point of RSVP fields and Polyxan hypergraphs,
\[
c_{\Polyx}[G_\ast] = c_{\RSVP}[\mathcal{L}_\ast].
\]
\end{theorem}

This is the first bridge between Parts~I and~III.

% ==============================================================
\section{Spectral Reconstruction Theorem}
% ==============================================================

\begin{theorem}[Spectral reconstruction]
A finite Polyxan hypergraph $G$ is uniquely determined, up to typed isomorphism,
by:
\begin{enumerate}
  \item the spectrum of $\Delta_G$, and  
  \item the sheaf cohomology groups $H^\bullet(G;\mathscr{F})$.
\end{enumerate}
\end{theorem}

\begin{proof}[Sketch]
The normalized spectrum determines all curvature-weighted incidence relations
up to typed symmetry.  
Sheaf cohomology resolves ambiguities in higher-order compatibility data.  
Because $\Polyx$ is equivalent to a presheaf topos, these invariants uniquely
determine the underlying object.
\end{proof}

Thus the pair $(\mathrm{Spec}(\Delta_G), H^\bullet(G;\mathscr{F}))$ constitutes a
semantic fingerprint of $G$.

% ==============================================================
\section{Summary}
% ==============================================================

Chapter~\ref{ch:spectral-semantics} established the spectral semantic theory of
Polyxan hypergraphs:

\begin{itemize}
  \item the normalized Laplacian $\Delta_G$ derived from the curvature-weighted
        incidence complex;
  \item the Polyxan heat kernel, heat trace, zeta function, and analytic torsion;
  \item the spectral action $S[G]=\mathrm{Tr}(\log\Delta_G)$, whose critical
        points are media quines;
  \item the Polyxan c-theorem: $c_{\Polyx}[G_t]$ strictly decreases under EHR;
  \item unification of $c_{\Polyx}$ with the RSVP central charge at the joint
        fixed point;
  \item the spectral reconstruction theorem, completing the invariant
        classification of Polyxan hypergraphs.
\end{itemize}

Spectral semantics forms the discrete harmonic backbone of Polyxan theory and
prepares the ground for the field–graph couplings and semantic galaxy formation
developed in Part~III.
